\clearpage\thispagestyle{empty}
\pdfbookmark[0]{Quellen zu Komet (25. Januar 1812)}{comet1812jan25-sources}
\begin{center}\textbf{Quellen zu Komet (25. Januar 1812)}\end{center}
\noindent Diese KI-Übertragung ins heutige Deutsch folgt den \href{https://gdz.sub.uni-goettingen.de/id/PPN236007467}{\textit{Werken}, Band~VI} (Göttingen, 1874), S.~345--347. Der Bericht erschien zuerst unter der Überschrift Göttingen in den \textit{Göttingischen gelehrten Anzeigen}, Heft~15 vom 25. Januar 1812, S.~\href{https://api.digitale-sammlungen.de/iiif/image/v2/bsb10538681_00191/full/1600,/0/default.jpg}{147}, \href{https://api.digitale-sammlungen.de/iiif/image/v2/bsb10538681_00192/full/1600,/0/default.jpg}{148}, \href{https://api.digitale-sammlungen.de/iiif/image/v2/bsb10538681_00193/full/1600,/0/default.jpg}{149} und \href{https://api.digitale-sammlungen.de/iiif/image/v2/bsb10538681_00194/full/1600,/0/default.jpg}{150}. Das digitalisierte Exemplar stammt aus der Bayerischen Staatsbibliothek.
\par\smallskip\noindent Die Überschrift der Werkausgabe nennt S.~147--149. Die am Ende des Berichts erwähnte fünfzehnzeilige Ephemeride steht in der Zeitschrift auf S.~150. In der Werkausgabe endet der Bericht mit dem Absatz, der die Ephemeride ankündigt. Die Schreibweise der Formeln und die Tabellenwerte folgen der Werkausgabe, einschließlich der geschwungenen Form der zweiten Konstanten. KI-Übertragung und KI-Satz.
\clearpage
